\documentclass[a4paper]{article}

% 文章排版
\usepackage{indentfirst}
\setlength{\parindent}{0em}
\usepackage{autobreak}
\allowdisplaybreaks
\usepackage{parskip}
\usepackage{geometry}
\geometry{top=3cm,bottom=3cm,left=2cm,right=2cm}
\usepackage{anyfontsize}

% 数学物理宏包
\usepackage{amsmath}
\usepackage{physics}
\renewcommand\div\divisionsymbol
\DeclareMathOperator{\diag}{diag}
\DeclareMathOperator{\sgn}{sgn}
\newcommand{\trans}{\text{\textsf{T}}}
\usepackage{tabularray}
\UseTblrLibrary{amsmath}

\usepackage{xcolor}
\usepackage{fontspec}
\setmainfont{STIX Two Text}
\usepackage{unicode-math}
\unimathsetup{mathrm=sym,mathbf=sym,mathsf=sym,bold-style=ISO}
\setmathfont{STIX Two Math}

\usepackage[fontset=none]{ctex}
\setCJKmainfont{FZShuSong-Z01}[BoldFont=Source Han Sans CN,ItalicFont=FZKai-Z03]
\setCJKsansfont{Source Han Sans CN}

% 符号重定义
\AtBeginDocument{
    \let\uglyepsilon\epsilon
    \renewcommand\epsilon\varepsilon
    \let\uglyleq\leq
    \renewcommand\leq\leqslant
    \let\uglygeq\geq
    \renewcommand\geq\geqslant
}


% 题目
\title{\textbf{广义相对论（天文方向）\ \ \ \ 作业八}}
\author{国家天文台 张植竣}
\date{2025年10月31日}

\begin{document}

\maketitle

\begin{enumerate}
    \item 普通标量场$\varphi$的协变能动张量为：
    \[
        T_{\mu\nu} = (\nabla_\mu \varphi)(\nabla_\nu \varphi) - \frac{1}{2} g_{\mu\nu} \nabla^\sigma \varphi \nabla_\sigma \varphi + g_{\mu\nu} V(\varphi)
    \]
    在Schwarzschild背景下，标量场只是半径$r$的函数。已知Schwarzschild度规为：
    \[
        \dd s^2 = -\qty(1 - \frac{2M}{r}) \dd t^2 + \qty(1 - \frac{2M}{r})^{-1} \dd r^2 + r^2 \dd \theta^2 + r^2 \sin^2\theta \dd \phi^2
    \]
    请计算协变能动张量在该度规下的各个分量。

    \textcolor{blue}{
    在史瓦西度规下，标量场$\varphi$只是$r$的函数，则：
    \[
        \nabla_\mu \varphi = \pdv{\varphi}{r} \delta^r_{\mu}
    \]
    记$\dfrac{\dd\varphi}{\dd r} = \varphi'$，则协变能动张量的第一项为：
    \[
        A_{rr} = (\nabla_r \varphi)(\nabla_r \varphi) = \varphi'^2,\quad A_{tt} = A_{\theta\theta} = A_{\varphi\varphi} = 0
    \]
    同时：
    \begin{align*}
        \nabla^\sigma \varphi \nabla_\sigma \varphi
        & = g^{\mu\nu} (\nabla_\mu \varphi)(\nabla_\nu \varphi) \\
        & = g^{rr} (\nabla_r \varphi)(\nabla_r \varphi) \\
        & = \qty(1 - \frac{2M}{r}) \varphi'^2
    \end{align*}
    因此协变能动张量的第二项为：
    \[
        B_{\mu\nu} = -\frac{1}{2} g_{\mu\nu} \nabla^\sigma \varphi \nabla_\sigma \varphi = \frac{1}{2} g_{\mu\nu} \qty(1 - \frac{2M}{r}) \varphi'^2
    \]
    即：
    \[
        B_{tt} = \frac{1}{2} \qty(1 - \frac{2M}{r}) \varphi'^2,\quad B_{rr} = -\frac{1}{2} \varphi'^2
    \]
    \[
        B_{\theta\theta} = -\frac{1}{2} r^2 \qty(1 - \frac{2M}{r}) \varphi'^2,\quad B_{\varphi\varphi} = -\frac{1}{2} r^2 \sin^2\theta \qty(1 - \frac{2M}{r}) \varphi'^2
    \]
    第三项为$C_{\mu\nu} = g_{\mu\nu} V(\varphi)$，即：
    \[
        C_{tt} = -\qty(1 - \frac{2M}{r}) V(\varphi),\quad C_{rr} = \qty(1 - \frac{2M}{r})^{-1} V(\varphi),
    \]
    \[
        C_{\theta\theta} = r^2 V(\varphi),\quad C_{\varphi\varphi} = r^2 \sin^2\theta V(\varphi)
    \]
    最终各分量为：
    \[
        T_{tt} = \frac{1}{2} \qty(1 - \frac{2M}{r}) \varphi'^2 - \qty(1 - \frac{2M}{r}) V(\varphi)
    \]
    \[
        T_{rr} = \frac{1}{2} \varphi'^2 + \qty(1 - \frac{2M}{r})^{-1} V(\varphi)
    \]
    \[
        T_{\theta\theta} = -\frac{1}{2} r^2 \qty(1 - \frac{2M}{r}) \varphi'^2 + r^2 V(\varphi)
    \]
    \[
        T_{\phi\phi} = -\frac{1}{2} r^2 \sin^2\theta \qty(1 - \frac{2M}{r}) \varphi'^2 + r^2 \sin^2\theta V(\varphi)
    \]
    其余非对角项均为零。
    }

    \item 在FRW宇宙背景下，标量场只是时间$t$的函数，计算协变能动张量在FRW宇宙背景中的各个分量。已知FRW度规为：
    \[
        \dd s^2 = -\dd t^2 + a^2(t) \qty(\dd r^2 + r^2 \dd \theta^2 + r^2 \sin^2\theta \dd \phi^2)
    \]

    \textcolor{blue}{
    在FRW宇宙背景下，标量场$\varphi$只是时间$t$的函数，则：
    \[
        \nabla_\mu \varphi = \pdv{\varphi}{t} \delta^t_{\mu}
    \]
    记$\dfrac{\dd\varphi}{\dd r} = \dot{\varphi}$，则协变能动张量的第一项为：
    \[
        A_{tt} = (\nabla_t \varphi)(\nabla_t \varphi) = \dot{\varphi}^2,\quad A_{rr} = A_{\theta\theta} = A_{\varphi\varphi} = 0
    \]
    同时：
    \begin{align*}
        \nabla^\sigma \varphi \nabla_\sigma \varphi
        & = g^{\mu\nu} (\nabla_\mu \varphi)(\nabla_\nu \varphi) \\
        & = g^{tt} (\nabla_t \varphi)(\nabla_t \varphi) \\
        & = -\dot{\varphi}^2
    \end{align*}
    因此协变能动张量的第二项为：
    \[
        B_{\mu\nu} = -\frac{1}{2} g_{\mu\nu} \nabla^\sigma \varphi \nabla_\sigma \varphi = \frac{1}{2} g_{\mu\nu} \dot{\varphi}^2
    \]
    即：
    \[
        B_{tt} = -\frac{1}{2} \dot{\varphi}^2,\quad B_{rr} = \frac{1}{2} a^2(t) \dot{\varphi}^2,
    \]
    \[
        B_{\theta\theta} = \frac{1}{2} a^2(t) r^2 \dot{\varphi}^2,\quad B_{\varphi\varphi} = \frac{1}{2} a^2(t) r^2 \sin^2\theta \dot{\varphi}^2
    \]
    第三项为$C_{\mu\nu} = g_{\mu\nu} V(\varphi)$，即：
    \[
        C_{tt} = -V(\varphi),\quad C_{rr} = a^2(t) V(\varphi),
    \]
    \[
        C_{\theta\theta} = a^2(t) r^2 V(\varphi),\quad C_{\varphi\varphi} = a^2(t) r^2 \sin^2\theta V(\varphi)
    \]
    最终各分量为：
    \[
        T_{tt} = \frac{1}{2} \dot{\varphi}^2 - V(\varphi)
    \]
    \[
        T_{rr} = a^2(t) \qty[\frac{1}{2} \dot{\varphi}^2 + V(\varphi)]
    \]
    \[
        T_{\theta\theta} = a^2(t) r^2 \qty[\frac{1}{2} \dot{\varphi}^2 + V(\varphi)]
    \]
    \[
        T_{\phi\phi} = a^2(t) r^2 \sin^2\theta \qty[\frac{1}{2} \dot{\varphi}^2 + V(\varphi)]
    \]
    其余非对角项均为零。
    }

\end{enumerate}

\end{document}